\documentclass{article}

\usepackage[utf8]{inputenc}
\usepackage[T1]{fontenc}
\usepackage[english]{babel}
\usepackage{microtype}

\usepackage{amsmath, amssymb, amsthm}
\usepackage{mathtools}
\usepackage{totcount}
\usepackage{enumitem}

\usepackage{newtxtext,newtxmath}

\usepackage{titlesec}
\titleformat{\chapter}[display]
  {\normalfont\huge\bfseries} % Format of the text
  {Lecture \thechapter}       % Label (e.g., "Lecture 1")
  {0.5em}                     % Spacing between label and title
  {\Huge}                     % Format of the title

% Theorem-like environments
\newtheorem{theorem}{Theorem}[section]
\newtheorem{lemma}[theorem]{Lemma}
\newtheorem{proposition}[theorem]{Proposition}
\newtheorem{corollary}[theorem]{Corollary}

\theoremstyle{definition}
\newtheorem{definition}[theorem]{Definition}
\newtheorem{example}[theorem]{Example}

\theoremstyle{remark}
\newtheorem{remark}[theorem]{Remark}

\title{Practice Midterm 1 for Elementary Number Theory}
\author{Egbert Rijke}
\date{\today}

\usepackage{hyperref}
\hypersetup{
    colorlinks=true,
    linkcolor=purple,
    filecolor=magenta,
    urlcolor=blue,
    citecolor=purple,
    pdftitle={Lecture notes on Number Theory},
    pdfpagemode=FullScreen,
}
\usepackage[capitalize,english]{cleveref}

  % We set up the exercise environment, which produces list environment in a new
% unnumbered subsection that also gets mentioned in the table of contents.

\newtotcounter{exercisecounter}
\newcommand{\exercise}{\item\stepcounter{exercisecounter}}

\newlist{exenum}{enumerate}{1}
\setlist[exenum]{noitemsep,label=\arabic*}
  %,ref=\thechapter.\arabic*}

\crefname{exenumi}{Exercise}{Exercises}
  
\newlist{subexenum}{enumerate}{1}
\setlist[subexenum]{noitemsep,label=(\alph*),ref=\theexenumi.\alph*}
\crefname{subexenumi}{Exercise}{Exercises}
  
\newenvironment{exercises}
{%
\subsection*{Exercises}%
\addcontentsline{toc}{section}{Exercises}%
\sectionmark{Exercises}%
\begin{exenum}}
{%
%\addtocounter{exercisecounter}{\theexenumi}%
\end{exenum}}

\begin{document}
\maketitle

\begin{exercises}
  \exercise (10 points) Define the Pell numbers $P_n$ recursively by
  \begin{equation*}
    P_0 := 0, \quad P_1 := 1, \qquad\text{and}\qquad P_{n+2}= 2P_{n+1}+P_n.
  \end{equation*}
  Prove by induction that
  \begin{equation*}
    \sum_{k=0}^{n} P_k^2 = \frac{P_nP_{n+1}}{2}.
  \end{equation*}
  \exercise (10 points) Use Euclid's algorithm to compute the following greatest common divisors, and express them as linear combinations via B\'ezout's identity:
  \begin{subexenum}
  \item $\gcd(1023,255)$
  \item $\gcd(378,273)$
  \item $\gcd(396,312)$
  \item $\gcd(623,439)$
  \item $\gcd(987,567)$
  \item $\gcd(8128,496)$
  \end{subexenum}
  Note: In the actual midterm we will choose numbers that are easier to compute with by hand, and there will not be so many. These are for practicing your pen-and-paper computations.
  \exercise (10 points) Show that for any integer $x$, the fraction
  \begin{equation*}
    \frac{7x+3}{5x+2}
  \end{equation*}
  is in lowest terms.
  \exercise (10 points) Show that for any natural number $n>0$ there is a unique pair of numbers $(m,r)$ such that the only square dividing $r$ is $1$, and
  \begin{equation*}
    n = m^2r.
  \end{equation*}
  This is called the \emph{square-free decomposition} of $n$.
\end{exercises}
\end{document}
